\documentclass[10pt,a4paper]{amsart}
\usepackage[margin=19mm]{geometry}
\usepackage{amsmath,amssymb,mathtools}
\usepackage[scr=rsfs,cal=euler]{mathalfa}
\usepackage{xfrac}
\usepackage[unicode,hidelinks,bookmarks=false]{hyperref}
\newcommand{\ie}{i.e.\ }
\newcommand{\cf}{cf.\ }
\newcommand{\resp}{respectively\ }
\newcommand{\Subsection}{\subsection}
\providecommand{\nobreakdash}{\nobreak\hbox{-}}
\setlength{\emergencystretch}{2em}
\multlinegap0pt
\usepackage{bm}
\usepackage{bbm}
\usepackage{dsfont}
\usepackage{xspace}
\usepackage{cancel}
\usepackage{mathtools}
\usepackage{amsthm}
\newtheorem{lemma}{Lemma}
\title[Starter-Iterator Neural Operator: A Unified Architecture for]{Starter-Iterator Neural Operator: A Unified Architecture for High-Fidelity Forward and Inverse PDE Problems}
\author{Kuilin Qin, Lianfang Wang, Xu Sun, Jiwei Jia, Yu Wang, Yong Wang, Yuping Duan}
\date{Source: arXiv:2606.18305v2 (CC BY 4.0)}
\begin{document}
\maketitle

\noindent\emph{Source and licence.} Starter-Iterator Neural Operator: A Unified Architecture for High-Fidelity Forward and Inverse PDE Problems. Version arXiv:2606.18305v2 (CC BY 4.0), distributed under the Creative Commons Attribution 4.0 International licence, as recorded in the arXiv licence metadata for this exact version and shown on the abstract page https://arxiv.org/abs/2606.18305v2. Full rights notice: licenses/arxiv-2606.18305-CC-BY-4.0.txt.\par\smallskip

\noindent\textit{Editorial scope.} Counted unit: the Lemma whose source label is \texttt{lem:univ\_iter} in the article (source0.tex lines 2031--2098, source label \texttt{lem:univ\_iter}), delivered together with the article's own complete original proof of it (source0.tex lines 2100--2342, one \texttt{proof} environment of 243 lines). No mathematics is written, repaired or re-derived: statement and proof are the article's own text line for line. Required context retained uncounted: eq:multi\_step\_learnable\_iterator (lines 683--693). Every definition, display and label of those lines is retained unchanged. Omitted from this package and not redistributed: 1--682, 694--2030, 2099--2099, 2343--3799. Nothing inside a retained excerpt is omitted or altered: every retained body is delivered exactly as the article's lines are written, so no omission receipt of any kind accompanies this package. Citations: the retained excerpts carry no citation command at all, so no bibliography entry of the article is transcribed and no reference is redistributed. Reference adaptation: none was needed and none was made; every reference inside the delivered unit resolves inside it, so no unresolved reference or citation is produced. Printed slips kept literal: no printed word, symbol, hypothesis or number of the counted statement or of its proof was altered. Layout and class: the article's own class is \texttt{elsarticle} with packages including amssymb, rotating, amsmath, geometry, graphicx, booktabs, hyperref, amsfonts, amsthm, mathrsfs, xcolor, textcomp, pifont, manyfoot; the wrapper is the lane's pinned \texttt{amsart} editorial wrapper at 10pt on A4, which restates the theorem-like environments the retained text needs and adds the article's own macro definitions that the retained text needs (none), with extra wrapper packages (bm, bbm, dsfont, xspace, cancel, mathtools). The delivered numbering is the unit's own. Assets: no retained span contains a figure environment, a graphics-inclusion command, a PDF-page inclusion, a table or a TikZ picture; no image file is copied into this package, the package declares no asset dependency and no image file is redistributed with it. Licence: version arXiv:2606.18305v2 of this article is distributed under the Creative Commons Attribution 4.0 International licence, as recorded in the arXiv licence metadata for this exact version and shown on the abstract page https://arxiv.org/abs/2606.18305v2. A full rights notice is delivered at licenses/arxiv-2606.18305-CC-BY-4.0.txt. Characters: every retained excerpt is pure ASCII, so the delivered engine, XeLaTeX with its default Latin Modern text font, reports no missing character.
		\begin{equation}
			\label{eq:multi_step_learnable_iterator}
			\mathcal I_\theta^{n+1}(\boldsymbol u;\boldsymbol f)
			:=
			\mathcal I_{\theta^{(n)}}\left(
			\mathcal I_\theta^n(\boldsymbol u;\boldsymbol f),
			\boldsymbol f
			\right),
			\qquad
			n=0,\ldots,N-1.
		\end{equation}

			\begin{lemma}[Compact-uniform approximation of the iterator]
				\label{lem:univ_iter}
				Let $\mathcal C\subset\mathcal X$ be compact,
				let
				\[
				\mathcal A\in\mathcal L(\mathcal Y,\mathcal X),
				\qquad
				\mathcal B\in\mathcal L(\mathcal X,\mathcal Y),
				\]
				and let
				$\boldsymbol u^{(0)}:\mathcal C\to\mathcal Y$
				be continuous. 
				Then, for every $\epsilon>0$ and every fixed
				$N\in\mathbb N^+$, there exist finite-rank bounded linear operators
				\[
				\mathcal A_{\theta^{(n)}}
				\in\mathcal L(\mathcal Y,\mathcal X),
				\qquad
				\mathcal B_{\theta^{(n)}}
				\in\mathcal L(\mathcal X,\mathcal Y),
				\qquad
				n=0,\ldots,N-1,
				\]
				which define the step-dependent learned iteration
				\[
				\begin{aligned}
					\boldsymbol u_\theta^{(n+1)}(\boldsymbol f)
					=
					\mathcal I_{\theta^{(n)}}\left(
					\boldsymbol u_\theta^{(n)}(\boldsymbol f),
					\boldsymbol f
					\right)
					:=
					\boldsymbol u_\theta^{(n)}(\boldsymbol f)
					+
					\mathcal B_{\theta^{(n)}}
					\left(
					\boldsymbol f
					-
					\mathcal A_{\theta^{(n)}}
					\boldsymbol u_\theta^{(n)}(\boldsymbol f)
					\right),
					\qquad
					n=0,\ldots,N-1,
				\end{aligned}
				\]
				with $\boldsymbol u_\theta^{(0)}(\boldsymbol f)
				= \boldsymbol u^{(0)}(\boldsymbol f)$, 
				such that
				\[
				\sup_{\boldsymbol f\in\mathcal C}
				\left\|
				\mathcal I^N\left(
				\boldsymbol u^{(0)}(\boldsymbol f);
				\boldsymbol f
				\right)
				-
				\mathcal I_\theta^N\left(
				\boldsymbol u^{(0)}(\boldsymbol f);
				\boldsymbol f
				\right)
				\right\|_{\mathcal Y}
				<\epsilon.
				\]
				Here, $\mathcal I^N$ denotes $N$ successive applications of the exact
				iteration rule and $\mathcal I_\theta^N$ denotes the ordered application of $N$
				iteration steps with iteration-dependent parameters, defined as Eq.~\eqref{eq:multi_step_learnable_iterator}.
			\end{lemma}

			\begin{proof}
				For each $\boldsymbol f\in\mathcal C$, let
				$\{\boldsymbol u^{(n)}(\boldsymbol f)\}_{n=0}^{N}$ denote the exact
				iteration generated by
				\[
				\boldsymbol u^{(n+1)}(\boldsymbol f)
				=
				\boldsymbol u^{(n)}(\boldsymbol f)
				+
				\mathcal B
				\left(
				\boldsymbol f
				-
				\mathcal A\boldsymbol u^{(n)}(\boldsymbol f)
				\right),
				\qquad
				n=0,\ldots,N-1,
				\]
				with the same initial state
				$\boldsymbol u^{(0)}(\boldsymbol f)$ as the learned iterator.
				
				Since $\boldsymbol u^{(0)}$ is continuous and $\mathcal A$ and
				$\mathcal B$ are bounded linear operators, the map $\boldsymbol f
				\longmapsto\boldsymbol u^{(n)}(\boldsymbol f)$ is continuous for every fixed $n$. Because $\mathcal C$ is compact and
				$N$ is finite, the trajectory and residual sets
				\[
				\mathcal U_N
				:=
				\bigcup_{n=0}^{N-1}
				\left\{
				\boldsymbol u^{(n)}(\boldsymbol f):
				\boldsymbol f\in\mathcal C
				\right\}
				\subset\mathcal Y, \qquad
				\mathcal R_N
				:=
				\bigcup_{n=0}^{N-1}
				\left\{
				\boldsymbol f
				-
				\mathcal A\boldsymbol u^{(n)}(\boldsymbol f):
				\boldsymbol f\in\mathcal C
				\right\}
				\subset\mathcal X
				\]
				are compact. 
				
				{ Since $\mathcal X$ and $\mathcal Y$ are separable Hilbert spaces,
					finite-rank orthogonal projections converge strongly to the identity and
					uniformly on compact subsets. By composing $\mathcal A$ and $\mathcal B$
					with such projections, for every $\delta>0$ and
					$n=0,\ldots,N-1$, we may choose finite-rank bounded linear operators
					$\mathcal A_{\theta^{(n)}}$ and
					$\mathcal B_{\theta^{(n)}}$ such that
					\[
					\sup_{\boldsymbol u\in\mathcal U_N}
					\left\|
					\left(
					\mathcal A-\mathcal A_{\theta^{(n)}}
					\right)\boldsymbol u
					\right\|_{\mathcal X}
					\leq\delta, \qquad
					\sup_{\boldsymbol r\in\mathcal R_N}
					\left\|
					\left(
					\mathcal B-\mathcal B_{\theta^{(n)}}
					\right)\boldsymbol r
					\right\|_{\mathcal Y}
					\leq\delta.
					\]
					These are compact-uniform approximation bounds rather than global
					operator-norm approximation bounds. }
				Because the orthogonal projections used in this construction have norm at
				most one, the approximating operators may also be chosen such that
				\[
				\left\|
				\mathcal A_{\theta^{(n)}}
				\right\|_{\mathcal L(\mathcal Y,\mathcal X)}
				\leq \left\|
				\mathcal A
				\right\|_{\mathcal L(\mathcal Y,\mathcal X)},
				\qquad
				\left\|
				\mathcal B_{\theta^{(n)}}
				\right\|_{\mathcal L(\mathcal X,\mathcal Y)}
				\leq \left\|
				\mathcal B
				\right\|_{\mathcal L(\mathcal X,\mathcal Y)},
				\]
				Set $\Lambda
				:=
				\max
				\left\{
				\|\mathcal A\|_{\mathcal L(\mathcal Y,\mathcal X)},
				\|\mathcal B\|_{\mathcal L(\mathcal X,\mathcal Y)}
				\right\}$, 
				then
				\[
				\left\|
				\operatorname{Id}_{\mathcal Y}
				-
				\mathcal B_{\theta^{(n)}}
				\mathcal A_{\theta^{(n)}}
				\right\|_{\mathcal L(\mathcal Y)}
				\leq
				1+\Lambda^2
				=:M.
				\]
				Here the constants $\Lambda$ and $M$ are both independent of
				both $n$ and $\delta$.
				For $\boldsymbol f\in\mathcal C$, we define
				\[
				\boldsymbol e^{(n)}(\boldsymbol f)
				:=
				\boldsymbol u^{(n)}(\boldsymbol f)
				-
				\boldsymbol u_\theta^{(n)}(\boldsymbol f).
				\]
				Since the exact and learned iterations start from the same initial state, i.e. $\boldsymbol e^{(0)}(\boldsymbol f)=\boldsymbol 0$.
				Subtracting the two iteration rules gives
				\[
				\begin{aligned}
					\boldsymbol e^{(n+1)}(\boldsymbol f)
					={}&
					\left(
					\operatorname{Id}_{\mathcal Y}
					-
					\mathcal B_{\theta^{(n)}}
					\mathcal A_{\theta^{(n)}}
					\right)
					\boldsymbol e^{(n)}(\boldsymbol f)+
					\left(
					\mathcal B-\mathcal B_{\theta^{(n)}}
					\right)
					\big(
					\boldsymbol f
					-
					\mathcal A\boldsymbol u^{(n)}(\boldsymbol f)
					\big)+
					\mathcal B_{\theta^{(n)}}
					\left(
					\mathcal A_{\theta^{(n)}}-\mathcal A
					\right)
					\boldsymbol u^{(n)}(\boldsymbol f).
				\end{aligned}
				\]
				Taking the norm in $\mathcal Y$ yields
				\[
				\begin{aligned}
					\big\|
					\boldsymbol e^{(n+1)}(\boldsymbol f)
					\big\|_{\mathcal Y}
					\leq &
					\big\|
					\operatorname{Id}_{\mathcal Y}
					-
					\mathcal B_{\theta^{(n)}}
					\mathcal A_{\theta^{(n)}}
					\big\|_{\mathcal L(\mathcal Y)}
					\big\|
					\boldsymbol e^{(n)}(\boldsymbol f)
					\big\|_{\mathcal Y}
					+
					\big\|
					\left(
					\mathcal B-\mathcal B_{\theta^{(n)}}
					\right)
					\big(
					\boldsymbol f
					-
					\mathcal A\boldsymbol u^{(n)}(\boldsymbol f)
					\big)
					\big\|_{\mathcal Y}
					\\
					&+
					\big\|
					\mathcal B_{\theta^{(n)}}
					\big\|_{\mathcal L(\mathcal X,\mathcal Y)}
					\big\|
					\left(
					\mathcal A_{\theta^{(n)}}-\mathcal A
					\right)
					\boldsymbol u^{(n)}(\boldsymbol f)
					\big\|_{\mathcal X}.
				\end{aligned}
				\]
				% Let $M:=1+C^2$.  
				% Then
				% \[
				% \left\|
				% \operatorname{Id}_{\mathcal Y}
				% -
				% \mathcal B_{\theta^{(n)}}
				% \mathcal A_{\theta^{(n)}}
				% \right\|_{\mathcal L(\mathcal Y)}
				% \leq M.
				% \]
				Using the compact-uniform approximation bounds, we obtain
				\[
				\big\|
				\boldsymbol e^{(n+1)}(\boldsymbol f)
				\big\|_{\mathcal Y}
				\leq
				M
				\big\|
				\boldsymbol e^{(n)}(\boldsymbol f)
				\big\|_{\mathcal Y}
				+
				(1+\Lambda)\delta.
				\]
				Since $\boldsymbol e^{(0)}(\boldsymbol f)=\boldsymbol 0$, recursive application gives
				\[
				\big\|
				\boldsymbol e^{(N)}(\boldsymbol f)
				\big\|_{\mathcal Y}
				\leq
				(1+\Lambda)\delta
				\sum_{k=0}^{N-1}M^k.
				\]
				The bound is uniform over $\boldsymbol f\in\mathcal C$. Since $N$ is fixed, we may choose $\delta>0$ sufficiently small such that
				\[
				(1+\Lambda)\delta
				\sum_{k=0}^{N-1}M^k
				<\epsilon.
				\]
				Consequently,
				\[
				\sup_{\boldsymbol f\in\mathcal C}
				\Big\|
				\mathcal I^N\left(
				\boldsymbol u^{(0)}(\boldsymbol f);
				\boldsymbol f
				\right)
				-
				\mathcal I_\theta^N\big(
				\boldsymbol u^{(0)}(\boldsymbol f);
				\boldsymbol f
				\big)
				\Big\|_{\mathcal Y}
				<\epsilon,
				\]
				which completes the proof.
			\end{proof}

\end{document}
